\section{Conventions and standard finite-action lemmas}
\label{app:finite-action}

This appendix isolates the elementary facts used to convert the finite group
action into signature and hitting-set statements.  They are included so that
the computational objects can be read against explicit mathematical
semantics.

\begin{lemma}[Orbit-sum invariance and Boolean evaluation]
For $T\subseteq\Omega_n$ and $g\in G$,
\[
 q^G_{[T]}(g\cdot P)=q^G_{[T]}(P).
\]
On Boolean patterns, this common value is the number of supports in
$G\cdot T$ contained in $P$.
\end{lemma}

\begin{proof}
Relabelling permutes the support orbit:
\[
\begin{aligned}
q^G_{[T]}(g\cdot x)
 &=\sum_{U\in G\cdot T}m_U(g\cdot x)
  =\sum_{U\in G\cdot T}m_{g^{-1}U}(x)\\
 &=\sum_{W\in G\cdot T}m_W(x).
\end{aligned}
\]
For a Boolean pattern, $m_U(P)$ is one exactly when $U\subseteq P$ and is zero
otherwise.
\end{proof}

\begin{lemma}[Orbit-incidence normalization]
With $A_{R,C}$ and $M_{R,C}$ as in Section~\ref{sec:model},
\[
 |R|A_{R,C}=|C|M_{R,C}.
\]
\end{lemma}

\begin{proof}
Count the set
\[
 I_{R,C}=\{(P,T)\in R\times C:T\subseteq P\}
\]
by $P$ and then by $T$.  Transitivity makes the two fibre sizes constant,
giving $|I_{R,C}|=|R|A_{R,C}=|C|M_{R,C}$.
\end{proof}

\begin{lemma}[Pair-separation equivalence]
For a finite dictionary $\Dict$ and subfamily $F\subseteq\Dict$, the restricted
signature $\Fingerprint_F$ is injective on pattern orbits if and only if
$F\cap\Sep_{\Dict}(R,R')\neq\varnothing$ for every pair of distinct
pattern orbits $R,R'$.  Empty separator masks must be included: if one
occurs, no subfamily of $\Dict$ can separate all pattern orbits.
\end{lemma}

\begin{proof}
The signatures of $R$ and $R'$ agree on $F$ precisely when every coordinate in
$F$ has equal values on the two rows, which is equivalent to
$F\cap\Sep_{\Dict}(R,R')=\varnothing$.
\end{proof}

\begin{lemma}[Duplicate and antichain reductions]
Let $\mathcal S$ be a family of nonempty subsets of a finite coordinate set.
Deleting a duplicate set, or deleting $T\in\mathcal S$ when some
$S\in\mathcal S$ satisfies $S\subsetneq T$, does not change the family of
hitting sets.
\end{lemma}

\begin{proof}
A duplicate imposes the same condition twice.  If a family hits $S$, then it
also hits every superset $T$; hence the latter condition is redundant.
\end{proof}

\begin{lemma}[Canonical representatives preserve the finite model]
Choosing the minimum unsigned bit mask in every pattern and support orbit
induces bijections between abstract orbits and registry rows.  Evaluation on
one canonical pattern representative therefore determines the coordinate on
the entire pattern orbit.
\end{lemma}

\begin{proof}
Every finite orbit has a unique minimum mask.  Distinct orbits are disjoint and
hence have distinct minima.  Orbit-sum invariance then gives the evaluation
claim.
\end{proof}

The last lemma explains why comparing canonical registries and signature
matrices is stronger than comparing only aggregate counts: it aligns every
abstract orbit and every coordinate before row-level comparison.
